% Part: history
% Chapter: biographies 
% Section: kurt-goedel

\documentclass[../../../include/open-logic-section]{subfiles}

\begin{document}

\olfileid[hi]{his}{bio}{god}

\olsection{कुर्ट ग्योडेल}

\olphoto{goedel-kurt}{कुर्ट ग्योडेल}

कुर्ट ग्योडेल (\textsc{गर}-डल) का जन्म 28 अप्रैल 1906 को ऑस्ट्रो-हंगेरियाई
साम्राज्य के ब्र्युन (अब चेक गणराज्य का ब्रनो) में हुआ। जिज्ञासु और कुशाग्र
स्वभाव के कारण परिवार बालक कुर्टेले को प्रायः ``Der kleine Herr Warum''
(नन्हे श्रीमान क्यों) कहता था। प्राथमिक विद्यालय से ही वे पढ़ाई में उत्कृष्ट
थे; केवल गणित में उन्हें सर्वोच्च से कम अंक मिला। खराब स्वास्थ्य के कारण
ग्योडेल प्रायः विद्यालय से अनुपस्थित रहते और उन्हें शारीरिक शिक्षा से छूट
थी। बचपन में उन्हें आमवाती ज्वर का निदान हुआ। चिकित्सकीय आकलन के विपरीत होने
पर भी वे जीवन भर मानते रहे कि इससे उनके हृदय पर स्थायी प्रभाव पड़ा था।

ग्योडेल ने 1924 में वियना विश्वविद्यालय में अध्ययन आरंभ किया और~1929 में
डॉक्टरेट पूरा किया। पहले वे भौतिकी पढ़ना चाहते थे, किंतु दार्शनिक रूडोल्फ़
कार्नाप के प्रभाव के कारण उनकी रुचि शीघ्र ही गणित और विशेषतः तर्कशास्त्र की
ओर चली गई। हान्स हान के निर्देशन में लिखे उनके शोधप्रबंध ने तादात्म्य सहित
प्रथम-क्रम विधेय तर्क का पूर्णता प्रमेय सिद्ध किया \citep{Godel1929}। केवल
एक वर्ष बाद उन्होंने अपने सबसे प्रसिद्ध परिणाम---पहला और दूसरा अपूर्णता
प्रमेय---प्राप्त किए (\citealt{Godel1931} में प्रकाशित)। वियना में अपने समय
के दौरान ग्योडेल, वियना मंडल में बहुत सक्रिय थे। यह कार्नाप सहित वैज्ञानिक
दृष्टि वाले दार्शनिकों का समूह था, जिनका कार्य ग्योडेल के परिणामों से विशेष
रूप से प्रभावित हुआ।

1938 में ग्योडेल ने अडेल निम्बुर्स्की से विवाह किया। उनके माता-पिता प्रसन्न
नहीं थे: वह उनसे छह वर्ष बड़ी और पहले से तलाक़शुदा ही नहीं थीं, बल्कि एक
रात्रिक्लब में नर्तकी का काम भी करती थीं। किंतु सामाजिक दबाव ने ग्योडेल को
प्रभावित नहीं किया और उनकी मृत्यु तक दोनों का सुखी विवाह बना रहा।

1938 में नाज़ी जर्मनी द्वारा ऑस्ट्रिया का विलय किए जाने के बाद ग्योडेल और
अडेल संयुक्त राज्य अमेरिका चले गए, जहाँ उन्होंने न्यू जर्सी के प्रिंसटन में
इंस्टीट्यूट फॉर एडवांस्ड स्टडी में पद ग्रहण किया। अंतर्मुखता और विलक्षण
स्वभाव के बावजूद प्रिंसटन में ग्योडेल का समय सहयोगपूर्ण और फलदायी रहा।
उन्होंने समुच्चय सिद्धांत, दर्शन और भौतिकी में निबंध प्रकाशित किए। विशेष रूप
से आईएएस में अपने सहकर्मी अल्बर्ट आइंस्टाइन से उनकी बहुत गहरी मित्रता हुई।

अंतिम वर्षों में ग्योडेल का मानसिक स्वास्थ्य बिगड़ गया। 1977 में उनकी पत्नी
के अस्पताल में भर्ती होने के कारण वह उनके लिए भोजन नहीं बना सकीं। जीवन भर
मानसिक स्वास्थ्य समस्याओं से जूझे ग्योडेल अंततः उत्पीड़न-भ्रम से ग्रस्त हो
गए। विष दिए जाने के घोर भय से ग्योडेल ने भोजन करने से इनकार कर दिया। 14
जनवरी 1978 को प्रिंसटन में भुखमरी से उनका निधन हुआ।

\begin{reading}
ग्योडेल के जीवन की पूर्ण जीवनी के लिए \citet{Dawson1997} देखें। अन्य
जीवनीपरक लेखों और तर्कशास्त्र तथा दर्शन में ग्योडेल के योगदान पर निबंधों के
लिए \citet{Wang1990}, \citet{Baaz2011}, \citet{Takeuti2003} और
\citet{Sigmund2007} देखें।

ग्योडेल का पीएचडी शोधप्रबंध मूल जर्मन में उपलब्ध है \citep{Godel1929}।
अपूर्णता प्रमेयों का मूल पाठ \citep{Godel1931} है। पत्राचार के एक चयन सहित
ग्योडेल के सभी प्रकाशित और अप्रकाशित लेख उनके \emph{Collected Papers}
\citet{Godel1986,Godel1990} में अंग्रेज़ी में उपलब्ध हैं।

ग्योडेल के अपूर्णता प्रमेयों के विस्तृत विवेचन के लिए \citet{Smith2013}
देखें। ग्योडेल के प्रमेयों की अनौपचारिक दार्शनिक चर्चा के लिए मार्क
लिन्सेनमायर का पॉडकास्ट \citep{Linsenmayer2014} देखें।
\end{reading}

\end{document}
